% Figure: geometric convergence of power iteration on a log-y axis,
% for three spectral-gap ratios. The straight lines confirm the
% error decays like (lambda_2/lambda_1)^t.
\begin{figure}[htbp]
\centering
\begin{tikzpicture}
\begin{semilogyaxis}[
  width=13cm,
  height=8.5cm,
  xlabel={iteration $t$},
  ylabel={eigenvalue error $\lvert \lambda_{1}^{(t)} - \lambda_{1}\rvert$},
  xmin=0, xmax=70,
  ymin=1e-7, ymax=2,
  tick label style={font=\footnotesize},
  label style={font=\small},
  legend style={font=\footnotesize, at={(0.97,0.97)}, anchor=north east},
  title={Power iteration: geometric convergence at three spectral gaps},
  title style={font=\bfseries},
  ymajorgrids=true,
  grid style={gray!25},
]
% r = 0.5: fast.
\addplot[engblue, very thick, domain=0:24, samples=25] {0.5^x};
\addlegendentry{$r = |\lambda_{2}/\lambda_{1}| = 0.5$}
% r = 0.7: medium.
\addplot[engamber, very thick, domain=0:45, samples=46] {0.7^x};
\addlegendentry{$r = 0.7$}
% r = 0.9: slow.
\addplot[engred, very thick, domain=0:70, samples=71] {0.9^x};
\addlegendentry{$r = 0.9$}
% Reference: three-digit accuracy line.
\addplot[gray, thin, dashed, domain=0:70] {1e-3};
\node[font=\footnotesize, gray, anchor=south west] at (axis cs:1,1.3e-3)
  {three-digit accuracy};
\end{semilogyaxis}
\end{tikzpicture}
\caption[Geometric convergence of power iteration]{The error of the
power-iteration eigenvalue estimate against iteration count, on a
logarithmic vertical axis, for three values of the spectral-gap ratio
$r = \lvert\lambda_{2}/\lambda_{1}\rvert$. The straight lines are the
signature of geometric convergence: the error after $t$ steps is
proportional to $r^{t}$, so $\log(\text{error})$ falls linearly with
slope $\log r$. A healthy gap ($r = 0.5$, blue) reaches three-digit
accuracy in about ten steps; a marginal gap ($r = 0.9$, red) needs
roughly sixty-six. The penalty for a closing gap is severe and
motivates the shift-and-invert and Krylov methods that widen the
effective gap. The Rayleigh-quotient eigenvalue estimate plotted here
converges at rate $r^{2t}$ for a symmetric matrix, twice the slope of
the eigenvector error, because the Rayleigh quotient is stationary at
the eigenvectors.}
\label{fig:vol02:ch10:power-convergence}
\end{figure}
